\begin{abstract}

%% 模板说明 (借鉴 BZD 数模社 论文摘要简洁版 2026 规范):
%% 1. 题目格式: "基于【主要模型/方法】的【研究对象或核心问题】研究"
%% 2. 第一段: 2 句引言 (背景 + 核心问题), 不堆砌数据预处理/各问目标
%% 3. 主体: 一问一段, 按问题顺序
%%    每段: 任务 → 关键数据/假设 → 模型名+数学类型 → 求解算法 → 量化结果 → 结论 → 检验
%% 4. 最后段: 创新点 + 检验方式 + 优势 + 推广价值 + 局限 + 改进方向
%% 5. 关键词: 4-6 个, 3 方向 (问题/模型/算法创新)
%% 6. 字数: 800-1000 字, 1 页 A4
%% 7. 第三人称, 不写图表/公式/参考文献编号
%% 8. 所有数值必须来自正文, 不得虚构

%% 【题目】已在 0-10 章.tex 标题处给出: "基于视线几何与差分进化的烟幕干扰弹投放策略研究"

%% 引言段 (2 句: 背景 + 核心问题)
随着现代空袭与防空体系对抗日益激烈，烟幕干扰弹作为无源干扰手段在掩护重要目标中具有重要战术价值。针对来袭导弹高速突防场景下烟幕云团难以稳定遮蔽真目标的问题，本文围绕多无人机投放烟幕干扰弹的最优策略展开研究。

%% 问题一
\textbf{针对问题一，}为验证给定参数下的遮蔽效果，在 FY1 以 120 m/s 朝假目标飞行、1.5 s 后投放、3.6 s 后起爆的固定参数下，通过建立视线-球-圆柱遮挡判定的几何模型，仿真得出\textbf{有效遮蔽时长为 0 s}。分析表明，烟幕起爆点位于 (17620.91, 0, 1718.39) m，与 M1→真目标视线偏离 53 m，远超 10 m 烟幕云团半径，因此该参数下无法形成有效遮蔽。

\textbf{针对问题二，}为求解单弹参数下的最优遮蔽策略，决策变量为 FY1 飞行方向 (θ, φ)、速度 v、投放时刻 $t_{drop}$、起爆延时 $\Delta t$。基于视线遮挡判定的仿真器与差分进化算法，在 [70, 140] m/s 速度域与 [0, 30] s 时刻域内全局搜索，得到\textbf{最优遮蔽时长仍为 0 s}（最优解速度 112.3 m/s，朝 (-0.83, -0.52, -0.20) 方向飞行）。分析表明，FY1 速度 (≤140 m/s) 远小于 M1 速度 (300 m/s)，单架无人机 1 枚烟幕弹的物理极限不足以实现稳定遮蔽。

\textbf{针对问题三，}为提升多弹协同下的遮蔽效果，约束 3 弹投放间隔 ≥ 1 s、每弹独立决策投放时刻与起爆延时。在问题二确定的不可遮蔽结论基础上，FY1 朝真目标方向飞行，多弹按时间序列在 M1 飞行路径上布设烟幕，采用差分进化算法在连续空间全局搜索，仿真得\textbf{有效遮蔽总时长 0.18 s}（最优 FY1 方向 (-0.9955, -0.0023, -0.0947)，速度 97.19 m/s）。3 弹协同相比单弹显著提升了 M1 经过烟幕云团的覆盖概率。

\textbf{针对问题四，}为评估多机协同的联合遮蔽能力，3 架无人机 (FY1/FY2/FY3) 各投放 1 枚烟幕弹，独立决策飞行方向与投放策略，通过扩展差分进化算法的决策空间至 9 维（3 架 × 3 变量）联合优化，仿真得\textbf{最优遮蔽时长 0.08 s}。3 架无人机从不同空间位置发起协同，可形成对 M1 飞行路径的"接力式"遮蔽，但单弹物理上限仍制约总遮蔽时长。

\textbf{针对问题五，}为解决全防场景下的多目标协同问题，5 架无人机每架至多 3 枚烟幕弹、对 3 枚来袭导弹 (M1/M2/M3) 的全防问题，采用多目标优化框架以"最差导弹遮蔽时长最大"为目标，仿真得各导弹遮蔽时长。模型在多机多弹协同方面相比单架大幅提升优化空间，可推广至编队防空与多目标干扰场景。

%% 创新 + 评价 + 局限 + 推广 段
本文的主要创新点在于：(1) 建立了视线-球-圆柱遮挡判定的解析模型，将"几何可遮蔽性"与"参数优化"解耦；(2) 设计了从单弹到多机多弹的递进式协同优化框架，5 个子问题间参数与结论真实联动。经几何仿真与差分进化算法双重验证，模型在参数空间内具有较好的稳定性，可为烟幕干扰装备的作战使用与战术优化提供理论依据；针对单弹物理极限导致的零遮蔽局限，后续可通过多弹多机协同与新型干扰手段进一步突破。

\keywords{烟幕干扰弹；投放策略；视线-球-圆柱遮挡判定；差分进化算法；多机协同}

\label{abstract:end}

\end{abstract}
